\def\level{Première}  % Classe à droite
\def\course{NSI}
\def\chaptername{Listes et Tuples} % Titre au centre
\def\docname{AP05}        % Identifiant à gauche

\pagestyle{cours_FP}
\makeheader{} % Affiche le bandeau

\section{Les listes}


\begin{definition}[Les listes]{}
	Une \textbf{liste} est une structure de donnée qui contient des éléments ou items identifiés par des indices.

Ces éléments sont:
\begin{itemize}
\item \textbf{ordonnés}: les éléments d'une liste sont rangés dans un ordre précis. Chaque élément a une position dans la liste, appelée indice.
\item \textbf{séparés par des virgules}: les éléments d'une liste sont séparés par des virgules.
\item \textbf{entre crochets}: les éléments d'une liste sont entourés de crochets \verb|[ ]| en python.
\end{itemize}


\medskip

\begin{tabular}{|p{2cm}|p{1cm}|p{1cm}|p{1cm}|p{1cm}|p{1cm}|p{1cm}|p{1cm}|p{1cm}|}
\hline
Indice \verb|i| & 0 & 1 & 2 & 3 & 4 & 5 & 6 & 7 \\ \hline
Elément & \verb|T[0]| & \verb|T[1]| & \verb|T[2]| & \verb|T[3]| & \verb|T[4]| & \verb|T[5]| & \verb|T[6]| & \verb|T[7]| \\ \hline

\end{tabular}



\end{definition}
\begin{easybox}{Créer une liste}{}
	\begin{minipage}{0.48\linewidth}
		\begin{CodePythontexAlt}[Largeur=1\linewidth,Centre,PremLigne=1, Lignes=false]{}
ma_liste = [1, 2, 3, 4, 5]
		\end{CodePythontexAlt}
	\end{minipage}\hfill
	\begin{minipage}{0.48\linewidth}
		\begin{CodePythontexAlt}[Largeur=1\linewidth,Centre,PremLigne=1, Lignes=false]{}
ma_liste = ['a', 'b', 'c']
		\end{CodePythontexAlt}
	\end{minipage}

	\smallskip

		\begin{ConsolePythontex}[Label=false]{}
ma_liste = [1, 2, 3, 4, 5]
print(type(ma_liste))		# Une liste est de type 'list'
		\end{ConsolePythontex}
\end{easybox}

\begin{easybox}{Accéder aux éléments d'une liste}{}
	Pour accéder à un élément d'une liste, on utilise son indice. Les indices commencent à 0.

	\smallskip

	\begin{minipage}{0.48\linewidth}
		\begin{ConsolePythontex}[Label=false]{}
ma_liste = [1, 2, 3, 4, 5]
print(ma_liste[0]) 
print(ma_liste[2]) 
		\end{ConsolePythontex}
	\end{minipage}\hfill
	\begin{minipage}{0.48\linewidth}
		\begin{ConsolePythontex}[Label=false]{}
ma_liste = [1, 2, 3, 4, 5]
print(ma_liste[-1]) 
print(ma_liste[-2])
		\end{ConsolePythontex}
	\end{minipage}
	\end{easybox}

	\begin{easybox}{Modifier les éléments d'une liste}{}
		Les éléments d'une liste peuvent être modifiés en utilisant leur indice.

	\begin{minipage}{0.48\linewidth}
		\begin{ConsolePythontex}[Label=false]{}
ma_liste = [1, 2, 3, 4, 5]
ma_liste[0] = 10
print(ma_liste)
		\end{ConsolePythontex}
	\end{minipage}\hfill
	\begin{minipage}{0.48\linewidth}
		\begin{ConsolePythontex}[Label=false]{}
ma_liste = [1, 2, 3, 4, 5]
ma_liste[4] = 20
print(ma_liste)
		\end{ConsolePythontex}
	\end{minipage}
\end{easybox}

\begin{easybox}{Connaitre le nombre d'éléments d'une liste}{}
	Pour connaitre le nombre d'éléments d'une liste, on utilise la fonction \verb|len()|.

	\begin{minipage}{0.48\linewidth}
		\begin{ConsolePythontex}[Label=false]{}
ma_liste = [1, 2, 3, 4, 5]
print(len(ma_liste))
		\end{ConsolePythontex}
	\end{minipage}\hfill
	\begin{minipage}{0.48\linewidth}
		\begin{ConsolePythontex}[Label=false]{}
ma_liste = ['a', 'b', 'c']
print(len(ma_liste))
		\end{ConsolePythontex}
	\end{minipage}
\end{easybox}

\begin{exercice_cours}[{Manipulation de base}]{}

	\begin{enumerate}
		\item Créer la liste suivante: \verb|[4, 8, 7, 9, 4]|
		\item Afficher le nombre d'éléments de cette liste. (\emph{affiche 5})
		\item Afficher le troisième élément de cette liste. (\emph{affiche 7})
		\item Modifier le deuxième élément de cette liste pour qu'il soit égal à 10
		\item Afficher la liste modifiée. (\emph{affiche [4, 10, 7, 9, 4]})
		\item Modifier le dernier élément de cette liste pour qu'il soit égal à 5
		\item Afficher la liste modifiée. (\emph{affiche [4, 10, 7, 9, 5]})
		\item Modifier le premier éléments de la liste en le remplaçant par la somme du premier et du troisième élément de cette liste
		\item Afficher la liste modifiée. (\emph{affiche [11, 10, 7, 9, 5]})

	\end{enumerate}
\end{exercice_cours}

\begin{easybox}{Itérer sur une liste avec une boucle FOR}{}

	Comme pour les chaines de caractères, il est possible d'itérer sur les éléments d'une liste à l'aide d'une boucle \verb|FOR|.	

	\smallskip
	\begin{minipage}{0.48\linewidth}
		\begin{CodePythontexAlt}[Largeur=1\linewidth,Centre,PremLigne=1]{}
ma_liste = [4, 8, 7, 9, 4]
for i in range(0, len(ma_liste)):
	print(ma_liste[i])
		\end{CodePythontexAlt}
	\end{minipage}\hfill
	\begin{minipage}{0.48\linewidth}
		\begin{CodePythontexAlt}[Largeur=1\linewidth,Centre,PremLigne=1]{}
ma_liste = [4, 8, 7, 9, 4]
for elt in ma_liste:
	print(elt)
		\end{CodePythontexAlt}
	\end{minipage}
\end{easybox}

\begin{exercice_cours}[Quelques algorithmes sur les listes]{}
	
	\begin{enumerate}
		\item Ecrire une fonction \texttt{somme(liste)} qui renvoie la somme des éléments d'une liste de nombres.
		
		\begin{CodePythontexAlt}[Largeur=1\linewidth,Centre,PremLigne=1, Lignes=false]{}
				def somme(liste):
					''' Renvoie la somme des éléments de liste'''







					#
				assert somme([1, 2, 3, 4, 5]) == 15
				assert somme([4, 8, 7, 9, 4]) == 32
				\end{CodePythontexAlt}

		\item Ecrire une fonction \texttt{maximum(liste)} qui renvoie le maximum des éléments d'une liste de nombres.
		\begin{CodePythontexAlt}[Largeur=1\linewidth,Centre,PremLigne=1, Lignes=false]{}
				def maximum(liste):
					''' Renvoie le maximum des éléments de liste'''









					#
				assert maximum([1, 2, 3, 4, 5]) == 5
				assert maximum([4, 8, 7, 9, 4]) == 9
				\end{CodePythontexAlt}
		\pagebreak

		\item Ecrire une fonction \texttt{est\_present(elt,liste)} qui renvoie \texttt{True} si l'élément \texttt{elt} est présent dans la liste \texttt{liste}, et \texttt{False} sinon.
		\begin{CodePythontexAlt}[Largeur=1\linewidth,Centre,PremLigne=1, Lignes=false]{}
				def est_present(elt, liste):
					''' Renvoie True si elt est dans liste et False sinon'''	









					#
				assert est_present(3, [1, 2, 3, 4, 5]) == True
				assert est_present(6, [1, 2, 3, 4, 5]) == False
				\end{CodePythontexAlt}
	
	\end{enumerate}
\end{exercice_cours}

\begin{easybox}{Liste de listes}{}
	Une liste peut contenir des éléments qui sont eux-mêmes des listes. On parle alors de liste de listes.

		\begin{ConsolePythontex}[Label=false]{}
ma_liste = [[1, 2, 3], [4, 5], [7, 8, 9, 10]]
print(len(ma_liste))
print(len(ma_liste[1]))
print(ma_liste[0])
print(ma_liste[1])
print(ma_liste[2])
print(ma_liste[0][1])
print(ma_liste[1][1])
print(ma_liste[2][0])
		\end{ConsolePythontex}

		

\end{easybox}

\begin{exercice_cours}[Manipulation de listes de listes]{}
	\begin{enumerate}
		\item Créer la liste suivante: \verb|[[1, 2], [3, 4], [5, 6]]|
		\item Afficher le nombre d'éléments de cette liste. (\emph{affiche 3})
		\item Afficher le deuxième élément de cette liste. (\emph{affiche [3, 4]})
		\item Afficher le premier élément du troisième élément de cette liste. (\emph{affiche 5})
		\item Modifier le deuxième élément du premier élément de cette liste pour qu'il soit égal à 10
		\item Afficher la liste modifiée. (\emph{affiche [[1, 10], [3, 4], [5, 6]]})
		\item Modifier le premier élément du deuxième élément de cette liste pour qu'il soit égal à 20
		\item Afficher la liste modifiée. (\emph{affiche [[1, 10], [20, 4], [5, 6]]})	
		\item Modifier le deuxième élément du troisième élément de cette liste en le remplaçant par la somme du premier et du deuxième élément de ce troisième élément
		\item Afficher la liste modifiée. (\emph{affiche [[1, 10], [20, 4], [5, 11]]})
	\end{enumerate}
\end{exercice_cours}

\pagebreak


\begin{exercice_cours}[Itérer sur une liste de listes]{}
	\begin{enumerate}
		\item Ecrire une fonction \texttt{somme\_liste\_de\_listes(liste)} qui renvoie la somme des éléments d'une liste de listes de nombres.
		
		\begin{CodePythontexAlt}[Largeur=1\linewidth,Centre,PremLigne=1, Lignes=false]{}
				def somme_liste_de_listes(liste):
					''' Renvoie la somme des éléments de liste'''









					#
				assert somme_liste_de_listes([[1, 2], [3, 4], [5, 6]]) == 21
				\end{CodePythontexAlt}

		\item Ecrire une fonction \texttt{maximum\_sous\_liste(liste)} qui renvoie la sous liste dont la somme des éléments est maximale.
		\begin{CodePythontexAlt}[Largeur=1\linewidth,Centre,PremLigne=1, Lignes=false]{}
				def maximum_sous_liste(liste):
					''' Renvoie la sous liste dont la somme des éléments est max'''










					#
				assert maximum_sous_liste([[3, 4, 16], [5, 6, 1, 8]]) == [3, 4, 16]
				\end{CodePythontexAlt}
	\end{enumerate}
\end{exercice_cours}

\begin{easybox}{Quelques méthodes sur les listes}{}

\begin{itemize}
\item \verb|append(elt)|: ajoute l'élément \verb|elt| à la fin de la liste	
\item \verb|pop()|: supprime et renvoie le dernier élément de la liste
\item \verb|pop(i)|: supprime et renvoie l'élément d'indice \verb|i| de la liste
\end{itemize}

\begin{ConsolePythontex}[Label=false]{}
ma_liste = [1, 2, 3, 4, 5]
ma_liste.append(6)
print(ma_liste)
a = ma_liste.pop()
print(a)
print(ma_liste)
b = ma_liste.pop(2)
print(b)
print(ma_liste)
		\end{ConsolePythontex}
\end{easybox}

\pagebreak

\begin{easybox}{Créer une liste par itération}{}

\begin{minipage}{0.48\linewidth}
		\begin{CodePythontexAlt}[Largeur=1\linewidth,Centre,PremLigne=1]{}
ma_liste = []
for i in range(0, 5):
	ma_liste.append(i)

print(ma_liste)		
# Affiche: [0, 1, 2, 3, 4]
		\end{CodePythontexAlt}
	\end{minipage}\hfill
	\begin{minipage}{0.48\linewidth}
		\begin{CodePythontexAlt}[Largeur=1\linewidth,Centre,PremLigne=1]{}
ma_liste = []
for i in range(0, 5):
	ma_liste = [i] + ma_liste

print(ma_liste)
# Affiche: [4, 3, 2, 1, 0]
		
		\end{CodePythontexAlt}
	\end{minipage}
\end{easybox}

\begin{easybox}{Créer une liste par compréhension}{}

	\begin{CodePythontexAlt}[Largeur=1\linewidth,Centre,PremLigne=1]{}
ma_liste = [i for i in range(0, 5)]
print(ma_liste)
# Affiche: [0, 1, 2, 3, 4]
		\end{CodePythontexAlt}
		\end{easybox}

\begin{exercice_cours}[Quelques méthodes sur les listes]{}

	\begin{enumerate}
		\item Ecrire une fonction \texttt{puissance\_2(n)} qui renvoie une liste contenant les puissances de 2 de 0 à n inclus.
		\begin{CodePythontexAlt}[Largeur=1\linewidth,Centre,PremLigne=1, Lignes=false]{}
				def puissance_2(n):
					''' Renvoie une liste contenant les puissances de 2 de 0 à n '''










					#
				assert puissance_2(3) == [1, 2, 4, 8]
				assert puissance_2(5) == [1, 2, 4, 8, 16, 32]
				\end{CodePythontexAlt}

		\item Ecrire une fonction \texttt{est\_pair(liste)} qui renvoie une liste contenant les éléments de liste qui sont pairs.
		\begin{CodePythontexAlt}[Largeur=1\linewidth,Centre,PremLigne=1, Lignes=false]{}
				def est_pair(liste):
					''' Renvoie une liste contenant les éléments pairs'''










					#
				assert est_pair([1, 2, 3, 4, 5]) == [2, 4]
				assert est_pair([1, 3, 5, 7, 9]) == []
				assert est_pair([2, 4, 6, 21, 8, 12, 13, 10]) == [2, 4, 6, 8, 12, 10]
				\end{CodePythontexAlt}

				\pagebreak

		\item Ecrire une fonction \texttt{supprime\_element(liste, element)} qui supprime les éléments d'une liste en partant du dernier jusqu'à l'obtention du premier \verb|0|.
		
		\begin{CodePythontexAlt}[Largeur=1\linewidth,Centre,PremLigne=1, Lignes=false]{}
				def supprime_element(liste):
					''' Supprime les éléments de liste en partant du dernier 
					jusqu'à l'obtention du premier 0'''









					#
				assert supprime_element([0, 0, 1, 1, 0, 1, 1, 1]) == [0, 0, 1, 1, 0]
				assert supprime_element([1, 1, 0, 1, 1, 1, 1, 1, 1]) == [1, 1, 0]
				\end{CodePythontexAlt}

\item Ecrire une fonction \texttt{est\_triee(liste)} qui renvoie \texttt{True} si la liste \texttt{liste} est triée dans l'ordre croissant, et \texttt{False} sinon.
					\begin{CodePythontexAlt}[Largeur=1\linewidth,Centre,PremLigne=1, Lignes=false]{}
					assert est_triee([1, 2, 3, 4, 5]) == True
					assert est_triee([5, 4, 3, 2, 1]) == False
					\end{CodePythontexAlt}
					
	\end{enumerate}
\end{exercice_cours}

\begin{easybox}{Le hasard et les listes}{}
	\begin{itemize}
		\item La fonction \verb|random.choice(liste)| permet de choisir un élément au hasard dans une liste.
		\item La fonction \verb|random.shuffle(liste)| permet de mélanger les éléments d'une liste de manière aléatoire.
	\end{itemize}

	\begin{ConsolePythontex}[Label=false]{}
import random
ma_liste = [1, 2, 3, 4, 5]
print(random.choice(ma_liste))
print(ma_liste)
random.shuffle(ma_liste)
print(ma_liste)
		\end{ConsolePythontex}
\end{easybox}

\begin{easybox}{Le slicing}{}
	Le slicing permet de créer une nouvelle liste à partir d'une liste existante en sélectionnant une portion de cette liste.

	\begin{CodePythontexAlt}[Largeur=1\linewidth,Centre,PremLigne=1]{}
ma_liste = [1, 2, 3, 4, 5]
print(ma_liste[1:4])
# Affiche: [2, 3, 4]
print(ma_liste[:3])
# Affiche: [1, 2, 3]
print(ma_liste[2:])
# Affiche: [3, 4, 5]
		\end{CodePythontexAlt}
\end{easybox}

\pagebreak

\begin{easybox}{Les effets de bord sur les listes}{}
	\begin{itemize}
		\item Les listes sont des objets mutables, c'est à dire que leur contenu peut être modifié après leur création.
		\item Lorsque l'on affecte une liste à une variable, on crée une référence à cette liste. Si on modifie la liste via cette variable, la modification sera visible via toutes les variables qui font référence à cette liste.
	\end{itemize}

	\begin{ConsolePythontex}[Label=false]{}
ma_liste = [1, 2, 3, 4, 5]
ma_liste2 = ma_liste
ma_liste2[0] = 10
print(ma_liste)
print(ma_liste2)
		\end{ConsolePythontex}
\end{easybox}

Si on veut faire une vraie copie d'une liste, on peut utiliser la syntaxe suivante:
	\begin{ConsolePythontex}[Label=false]{}
ma_liste = [1, 2, 3, 4, 5]
ma_liste2 = ma_liste[:]
ma_liste2[0] = 10
print(ma_liste)
print(ma_liste2)
		\end{ConsolePythontex}

Lorsqu'on utilise des listes dans des fonctions, il faut faire attention aux effets de bord. Si on modifie une liste passée en argument d'une fonction, la modification sera visible à l'extérieur de la fonction.

Dans cet exemple, il n'est pas utile de renvoyer la liste modifiée, car la modification est déjà visible à l'extérieur de la fonction.
\begin{CodePythontexAlt}[Largeur=1\linewidth,Centre,PremLigne=1]{}	
def ajoute_element(liste, element):
	liste.append(element)

ma_liste = [1, 2, 3, 4, 5]
ajoute_element(ma_liste, 6)
print(ma_liste)			# Affiche: [1, 2, 3, 4, 5, 6]
		\end{CodePythontexAlt}

			\begin{exercice_cours}[Effets de bord et fonctions]{}
	 Ecrire une fonction \texttt{decalage\_gauche(liste)} qui décale les éléments de la liste \texttt{liste} vers la gauche, en mettant le premier élément à la fin de la liste. La liste doit être modifiée en place.
					\begin{CodePythontexAlt}[Largeur=1\linewidth,Centre,PremLigne=1, Lignes=false]{}
					def decalage_gauche(liste):
						''' Décale les éléments de liste vers la gauche '''




						


					#	
					assert decalage_gauche([1, 2, 3, 4, 5]) == [2, 3, 4, 5, 1]
					assert decalage_gauche([5, 4, 3, 2, 1]) == [4, 3, 2, 1, 5]
					\end{CodePythontexAlt}

			\end{exercice_cours}

		\pagebreak

		\section{Les tuples}

		\begin{definition}[Les tuples]{}
		Un \textbf{tuple} est une structure de donnée qui contient des éléments ou items identifiés par des indices.	
		Les éléments d'un tuple sont:
\begin{itemize}
\item \textbf{ordonnés}: chaque élément a une position dans le tuple, appelée indice.
\item \textbf{séparés par des virgules}: les éléments d'un tuple sont séparés par des virgules.
\item \textbf{entre parenthèses}: les éléments d'un tuple sont entourés de parenthèses \verb|( )| en python.
\end{itemize}
\medskip
\begin{tabular}{|p{2cm}|p{1cm}|p{1cm}|p{1cm}|p{1cm}|p{1cm}|p{1cm}|p{1cm}|p{1cm}|}
\hline
Indice \verb|i| & 0 & 1 & 2 & 3 & 4 & 5 & 6 & 7 \\ \hline
Elément	& \verb|T[0]| & \verb|T[1]| & \verb|T[2]| & \verb|T[3]| & \verb|T[4]| & \verb|T[5]| & \verb|T[6]| & \verb|T[7]| \\ \hline

\end{tabular}

			
		\end{definition}

	
\begin{easybox}{Créer un tuple}{}
	\begin{minipage}{0.48\linewidth}
		\begin{CodePythontexAlt}[Largeur=1\linewidth,Centre,PremLigne=1, Lignes=false]{}
mon_tuple = (1, 2, 3, 4, 5)
		\end{CodePythontexAlt}
	\end{minipage}\hfill
	\begin{minipage}{0.48\linewidth}
		\begin{CodePythontexAlt}[Largeur=1\linewidth,Centre,PremLigne=1, Lignes=false]{}
mon_tuple = ('a', 'b', 'c')
		\end{CodePythontexAlt}
	\end{minipage}


		\begin{ConsolePythontex}[Label=false]{}
mon_tuple = (1, 2, 3, 4, 5)
print(type(mon_tuple))		# Un tuple est de type 'tuple'
		\end{ConsolePythontex}
\end{easybox}

\begin{easybox}{Accéder aux éléments d'un tuple}{}
	Pour accéder à un élément d'un tuple, on utilise son indice. Les indices commencent à 0.

	\smallskip

	\begin{minipage}{0.48\linewidth}
		\begin{ConsolePythontex}[Label=false]{}
mon_tuple = (1, 2, 3, 4, 5)
print(mon_tuple[0]) 
print(mon_tuple[2]) 
		\end{ConsolePythontex}
	\end{minipage}\hfill
	\begin{minipage}{0.48\linewidth}
		\begin{ConsolePythontex}[Label=false]{}
mon_tuple = (1, 2, 3, 4, 5)
print(mon_tuple[-1]) 
print(mon_tuple[-2])
		\end{ConsolePythontex}
	\end{minipage}
	\end{easybox}	

	\begin{propriete}[Immutabilité des tuples]{}
		Les éléments d'un tuple ne peuvent pas être modifiés après la création du tuple. Un tuple est une structure de donnée immuable.

	\end{propriete}

